\documentclass[conference]{IEEEtran}
\IEEEoverridecommandlockouts
% The preceding line is only needed to identify funding in the first footnote. If that is unneeded, please comment it out.
\usepackage{cite}
\usepackage{amsmath,amssymb,amsfonts}
\usepackage{algorithmic}
\usepackage{graphicx}
\usepackage{textcomp}
\usepackage{xcolor}
\usepackage{longtable}
\usepackage{booktabs}
\usepackage{array}
\usepackage{caption}
\usepackage{ragged2e}
\usepackage{svg}
\def\BibTeX{{\rm B\kern-.05em{\sc i\kern-.025em b}\kern-.08em
    T\kern-.1667em\lower.7ex\hbox{E}\kern-.125emX}}
\begin{document}

\title{A Quantum Threat Model for Hyperledger Fabric: Cryptographic Audit, Risk Prioritisation, and Formal
Integration Requirements\\
% {\footnotesize \textsuperscript{*}Note: Sub-titles are not captured in Xplore and
% should not be used}
% \thanks{Identify applicable funding agency here. If none, delete this.}
}

\author{\IEEEauthorblockN{Zara Laila Cheong}
\IEEEauthorblockA{\textit{School of Digital Science} \\
\textit{Univeristi Brunei Darussalam}\\
Bandar Seri Begawan, Brunei \\
zaralaila@icloud.com}
\and
\IEEEauthorblockN{Ali Tufail}
\IEEEauthorblockA{\textit{School of Digital Science} \\
\textit{Univeristi Brunei Darussalam}\\
Bandar Seri Begawan, Brunei \\
ali.tufail@ubd.edu.bn}
\and
\IEEEauthorblockN{Rosyzie Anna Apong}
\IEEEauthorblockA{\textit{School of Digital Science} \\
\textit{Univeristi Brunei Darussalam}\\
Bandar Seri Begawan, Brunei \\
rosyzie.apong@ubd.edu.bn}
}

\maketitle

\begin{abstract}
The advent of cryptographically relevant quantum computers threatens every cryptographic system whose security rests on the elliptic curve discrete logarithm problem, and Hyperledger Fabric, the dominant enterprise permissioned blockchain framework, depends uniformly on ECDSA and ECDH over the P-256 curve across its entire transaction lifecycle and identity infrastructure. This paper presents the first component-level quantum vulnerability audit and formal threat model for Hyperledger Fabric. Auditing version 2.5.12, we identify twenty-five distinct quantum-vulnerable cryptographic operations spanning five transaction lifecycle phases, the Fabric CA and Membership Service Provider identity subsystems, the gossip protocol, and the key management layer, and we confirm that the vulnerability is both complete and uniform: every public-key operation is simultaneously broken by Shor's algorithm. From this inventory we derive five formal threat scenarios, each with explicit adversarial capability assumptions, and we prioritise them in a risk matrix combining severity and exploitability. We show that the append-only immutable ledger amplifies the harvest-now-decrypt-later threat permanently, since every historical public key and signature remains accessible to a quantum adversary regardless of when migration is completed, and we adapt Mosca's inequality to this blockchain setting across five representative deployment scenarios. The analysis identifies the root CA self-signing operation and the orderer block-signing operation as the highest-priority migration targets, and it exposes a partial migration vulnerability in which post-quantum transaction signing provides no protection if the certificate authority remains classical. We conclude by deriving eight formal integration requirements, traceable to the identified threats, that constitute a complete specification for a quantum-resistant Fabric architecture, including the requirement that the live CA itself issue post-quantum certificates at runtime.
\end{abstract}

\begin{IEEEkeywords}
post-quantum cryptography, Hyperledger Fabric, permissioned blockchain, quantum threat model, harvest-now-decrypt-later, ECDSA, ML-DSA, cryptographic agility, security audit
\end{IEEEkeywords}

\section{Introduction}
\label{sec:tm_intro}

The emergence of quantum computing as a credible engineering target
introduces a class of cryptographic vulnerability that is qualitatively
different from the vulnerabilities addressed by conventional security
analysis. Classical security vulnerabilities typically arise from
implementation errors, misconfiguration, or the use of algorithms with
insufficient parameter sizes. The quantum vulnerability is structural:
every correctly implemented instance of a cryptographic algorithm family
whose security rests on the discrete logarithm problem or the integer
factorisation problem is rendered entirely insecure by Shor's
algorithm~\cite{shor1994} executing on a cryptographically relevant
quantum computer. No implementation improvement, no parameter adjustment,
and no operational change can restore security once the underlying
mathematical problem is tractable. Migration to post-quantum cryptographic
algorithms is the only effective countermeasure.

Hyperledger Fabric is the dominant enterprise permissioned blockchain
framework, deployed across finance, healthcare, supply chain, and
government applications~\cite{androulaki2018}. Its cryptographic
architecture relies uniformly on the Elliptic Curve Digital Signature
Algorithm with the P-256 curve for all digital signatures and on Elliptic
Curve Diffie-Hellman with the P-256 curve for all session key exchange
across its transaction lifecycle and identity management infrastructure.
Both primitives depend on the elliptic curve discrete logarithm problem,
which Shor's algorithm solves in polynomial time. A quantum adversary
capable of executing Shor's algorithm against P-256 can recover any
private key from its corresponding public key, forge any signature on
behalf of any identity, and decrypt any session key from a recorded TLS
handshake.

The harvest-now-decrypt-later threat is particularly severe for blockchain
deployments. Hyperledger Fabric's append-only immutable ledger permanently
preserves every transaction signature, every block header signature, and
every certificate chain ever committed. Standard key rotation procedures
do not remove this historical material from the ledger, and out-of-band
ledger reconfiguration is not part of normal Fabric operation; the
historical material therefore remains accessible to any entity with
ledger read access regardless of when a post-quantum migration is
completed. For Fabric networks
carrying data with long sensitivity lifetimes, the effective security
requirement extends not merely to future records but to the entire
accumulated history of the network.

Despite this urgency, no prior work has produced a component-level quantum
vulnerability audit for Hyperledger Fabric or for any enterprise
blockchain platform. Existing literature acknowledges the general quantum
threat to blockchain cryptography~\cite{fernandez2020, yang2024survey}
but does not enumerate specific vulnerable operations. Prior implementation
proposals proceed without formal threat analysis, leaving their designs
without systematic justification and unable to demonstrate completeness
with respect to the full attack surface~\cite{holcomb2021, campbell2019,
hu2025prototype, razdan2025postquantum}.

This chapter presents the following contributions. First, the first
component-level quantum vulnerability audit of Hyperledger Fabric,
identifying twenty-five distinct quantum-vulnerable cryptographic
operations across the full transaction lifecycle and identity management
infrastructure. Second, five formal threat scenarios derived from the
audit findings, each with explicit adversarial capability assumptions,
impact assessment, and harvest-now-decrypt-later dimensions specific to
the immutable ledger context. Third, a risk prioritisation matrix that
orders identified threats by severity and exploitability, with
architectural conclusions guiding migration priorities. Fourth, the
application of Mosca's inequality to five representative Fabric
deployment scenarios. Fifth, eight formal integration requirements, each
traceable to one or more identified threat scenarios, that constitute a
complete specification for quantum-resistant Fabric architecture and form
the direct input to Chapter~\ref{ch:architecture}.

The scope of the analysis is bounded by the standard Fabric deployment
model. Channel management operations that are identical in cryptographic
structure to transaction signing are included. Smart contract code
deployed on the chaincode runtime is excluded, as chaincode-level
cryptographic operations are application-layer choices orthogonal to
the infrastructure-layer vulnerabilities addressed here. External
hardware security module integrations are noted where relevant but not
analysed in detail, as they are deployment-specific configurations
rather than properties of the standard Fabric architecture.

The remainder of the chapter is organised as follows.
Section~\ref{sec:tm_method} describes the methodological framework and
audit scope. Section~\ref{sec:tm_inventory} presents the cryptographic
inventory of twenty-five operations. Section~\ref{sec:tm_threats}
constructs the five formal threat scenarios. Section~\ref{sec:tm_risk}
presents the risk prioritisation matrix. Section~\ref{sec:tm_mosca}
applies Mosca's inequality to representative deployment scenarios.
Section~\ref{sec:tm_requirements} derives the eight formal integration
requirements. Section~\ref{sec:tm_discussion} discusses the findings.
Section~\ref{sec:tm_conclusion} concludes.

\section{Methodology}
\label{sec:tm_method}

\subsection{Scope and Target System}

This chapter addresses Research Question RQ1: which cryptographic
components of Hyperledger Fabric are quantum-vulnerable, and what is
the severity and exploitability of each? The audit targets Hyperledger
Fabric v2.5.12, the most recent point release of the v2.5 long-term
support branch at the time of investigation. Scope encompasses all
components participating in the standard transaction lifecycle: the
client software development kit, endorsing peer nodes, the Raft ordering
service, Fabric CA, and the Membership Service Provider. Sources
consulted include the Fabric technical documentation, the BCCSP interface
specification, the \texttt{core.yaml} and \texttt{orderer.yaml}
configuration schemas, and the Fabric CA server configuration.

Two categories of operation fall outside scope. Hash functions including
SHA-256 and SHA-3 are noted as quantum-affected by Grover's algorithm,
which reduces the effective security of SHA-256 from 256 bits to 128
bits~\cite{grover1996}. Under current NIST guidance, 128 bits of
post-quantum security is considered acceptable, and SHA-256 is not
classified as requiring replacement~\cite{nist_sp80057}. AES-256,
used for private key storage at rest in the Fabric keystore, is similarly
affected by Grover's algorithm but retains 128 bits of post-quantum
security and is excluded from primary remediation scope. The audit
therefore concentrates on asymmetric cryptographic operations threatened
by Shor's algorithm.

\subsection{Analytical Framework}

Established vulnerability management frameworks were considered and
rejected as the primary analytical framework for this investigation.
The Common Vulnerabilities and Exposures classification focuses on
specific software implementation flaws rather than structural algorithmic
weaknesses. The OWASP Cryptographic Failures category addresses
implementation vulnerabilities such as weak random number generation,
incorrect cipher modes, and insufficient key lengths. NIST Special
Publication 800-57 provides key management recommendations designed for a
classical adversary model~\cite{nist_sp80057}.

None of these frameworks captures the structural quantum vulnerability
that motivates this investigation: the property that every correctly
implemented instance of ECDSA P-256 is simultaneously and completely
insecure against a quantum adversary, regardless of implementation quality
or key size within the standard parameter set. The framework adopted in
this chapter combines structured threat modelling following the approach
of Shostack~\cite{shostack2014}, drawing on the STRIDE classification
originally introduced by Kohnfelder and Garg in a 1999 Microsoft
internal memo for systematic enumeration of spoofing, tampering,
repudiation, information-disclosure, denial-of-service, and
elevation-of-privilege threats, and later popularised in the
\textit{TechNet Magazine} article cited here~\cite{stride2004},
with a quantum-specific vulnerability classification informed by the
NIST post-quantum standardisation literature~\cite{nist_ir8105, bernstein2017pqc}.

\subsection{Four-Stage Audit Process}

The audit proceeds through four stages. In the first stage, every major
architectural component of Fabric v2.5.12 is enumerated and the
cryptographic operations it performs within the standard transaction
lifecycle are identified. In the second stage, each operation is mapped
to its specific algorithm, key type, key size, and operational context.

In the third stage, each operation is classified with respect to quantum
vulnerability. Operations employing asymmetric algorithms based on the
elliptic curve discrete logarithm problem are classified as threatened by
Shor's algorithm. Operations employing symmetric algorithms or hash
functions are classified as Grover-affected. In the fourth stage, each
Shor-threatened operation is assessed for severity on a four-level scale
(Critical, High, Medium, Low) and for exploitability on a three-level
scale (High, Medium, Low).

\label{sec:tm_severity}
The severity scale is calibrated on the \textit{scope} of compromise
that successful exploitation enables, not on the \textit{reversibility}
of the resulting consequence. Reversibility is a property of every
committed operation in Fabric's lifecycle, since the ledger is uniformly
append-only; using it as the calibration criterion would collapse the
scale and make almost every Shor-threatened operation Critical, which
would defeat the purpose of prioritisation.

A \textbf{Critical} severity rating is assigned when exploitation enables
compromise of the integrity or authenticity of the entire Fabric network,
or of a trust anchor whose compromise propagates to all identities
certified under it. Operations rated Critical are those whose blast
radius is the network or an entire organisation, not a single
participant.

A \textbf{High} severity rating is assigned when exploitation enables
compromise of individual transactions, forged endorsements, or
impersonation of peer or client identities within a channel. Operations
rated High are those whose blast radius is bounded by a single recovered
identity and whose effect on the ledger is further constrained by the
applicable channel endorsement policy.

A \textbf{Medium} severity rating is assigned when exploitation enables
breaking the confidentiality of inter-node communications, including
decryption of previously recorded traffic.

A \textbf{Low} severity rating is assigned when exploitation yields only
partial information or marginal capability that does not directly enable
forgery or impersonation. The Low rating is defined here for
methodological completeness and to make the severity scale uniform; in
the present audit, no Shor-threatened operation in Hyperledger Fabric
v2.5.12 was rated Low, because every public-key operation in the
transaction lifecycle enables at least confidentiality compromise (when
the operation is an ECDH key exchange) or identity-bounded forgery
(when the operation is an ECDSA signature). Low severity remains in
the rubric so that the methodology is reusable for systems that may
contain genuinely marginal-impact ECDLP-dependent operations.



\begin{figure*}
    \centering
    \includesvg[width=1\linewidth]{figures/fig_ch4_methodology}
    \caption{Four-stage quantum vulnerability audit methodology for
Hyperledger Fabric v2.5.12}
    \label{fig:ch4_methodology}
\end{figure*}

\section{Cryptographic Inventory}
\label{sec:tm_inventory}

\subsection{Uniform ECDLP Dependence}

The cryptographic audit of Hyperledger Fabric v2.5.12 reveals a complete
and uniform dependence on the elliptic curve discrete logarithm problem
across all public-key cryptographic functions in the standard transaction
lifecycle. Every digital signature uses ECDSA with the NIST P-256 curve.
Every session key exchange uses ECDH with the NIST P-256 curve. There is
no cryptographic diversity in the public-key infrastructure: all asymmetric
security guarantees rest on the same mathematical hardness assumption and
are therefore simultaneously eliminated by Shor's algorithm.

The audit identifies twenty-five distinct quantum-vulnerable operations
organised across five transaction lifecycle phases, two identity management
subsystems, the gossip protocol, and the key management layer. These
operations are presented fully in Table~\ref{tab:crypto_inventory} and
discussed in the sections that follow.

The mathematical problem underlying this uniform dependence is the Elliptic
Curve Discrete Logarithm Problem. Given an elliptic curve $E$ over a prime
field $\mathbb{F}_p$, a base point $G \in E(\mathbb{F}_p)$ of prime order
$n$, and a point $Q \in E(\mathbb{F}_p)$, the ECDLP is to find the unique
integer $k \in \{1,\ldots,n-1\}$ such that

\begin{equation}
Q = kG.
\label{eq:ecdlp}
\end{equation}

The best known classical algorithm for solving~(\ref{eq:ecdlp}) is Pollard's
rho algorithm, requiring $O(\sqrt{n})$ group operations. For the P-256 curve,
$n \approx 2^{256}$, giving a classical complexity of approximately $2^{128}$
operations, considered computationally infeasible with foreseeable classical
hardware. Shor's algorithm reduces this to $O((\log n)^3)$ quantum gate
operations~\cite{shor1994}:

\begin{equation}
T_{\mathrm{classical}}(\mathrm{ECDLP}) = O\!\left(\sqrt{n}\right),
\qquad
T_{\mathrm{quantum}}(\mathrm{ECDLP}) = O\!\left((\log n)^3\right).
\label{eq:complexity}
\end{equation}

The consequence for ECDSA P-256 is categorical. Defining the security level
$\lambda$ as the binary logarithm of the number of operations required by the
best known attack, the transition from classical to quantum adversary is

\begin{equation}
\lambda_{\mathrm{classical}}(\mathrm{ECDSA\text{-}P256}) \approx 128\text{ bits},
\qquad
\lambda_{\mathrm{quantum}}(\mathrm{ECDSA\text{-}P256}) = 0.
\label{eq:security_bits}
\end{equation}

Equation~(\ref{eq:security_bits}) establishes that no implementation improvement,
parameter adjustment, or operational hardening can restore security once
Shor's algorithm is executable at the target key size. This is the fundamental
property that distinguishes the quantum vulnerability from all classical
vulnerabilities addressed by conventional security analysis frameworks.

\subsection{Phase 1: Transaction Proposal and Client Signing}

Op-01 is the client transaction proposal signing operation, in which
the client's ECDSA P-256 private key authenticates the client's identity
before submission to endorsing peers. The signature is constructed over
the proposal payload, the channel identifier, the chaincode specification,
and the transaction nonce, making it an authenticated binding of the
client's identity to a specific intended action. This signature is the
initial authentication gate for the transaction lifecycle; its forgery
enables a quantum adversary to submit arbitrary transaction proposals
under any client identity, including identities with elevated permissions
such as system chaincode invocation rights. Op-01 is rated High severity.

The severity of Op-01 is partly bounded by the endorsement policy
mechanism: a fraudulent transaction proposal signed under a compromised
client identity can only succeed in being endorsed if it satisfies the
chaincode's endorsement policy. However, in deployments where client
identities hold broad endorsement permissions, or where the endorsement
policy requires only a single organisation's signature, Op-01 compromise
can directly enable fraudulent state writes to the ledger. For chaincodes
that implement access-controlled operations such as asset transfer or
account balance modification, the impact of client identity compromise
at Op-01 is equivalent to unrestricted write access to the relevant ledger
state.

Op-02 is the client TLS handshake to the peer, which uses ECDH P-256 to
establish the session key for the proposal submission channel. The ECDH
exchange parameters are transmitted in cleartext in the TLS ClientHello
and ServerKeyExchange messages, making them available to any passive
observer of network traffic. Recorded handshakes are retroactively
decryptable by a quantum adversary who applies Shor's algorithm to the
recorded ECDH exchange, recovering the session key and decrypting all
application-layer traffic carried by that TLS session. Op-02 is rated
Medium severity, reflecting the confidentiality impact without direct
integrity consequences for the ledger.

\subsection{Phase 2: Peer Endorsement}

Op-03 is the verification of the client's ECDSA signature at the
endorsing peer. This operation is performed before the chaincode execution
begins; a client whose signature cannot be verified is rejected immediately
without committing any peer resources to chaincode simulation. Op-03
gates all transaction processing: its bypass allows fraudulent proposals
to enter the chaincode execution pipeline regardless of whether the
submitting client holds valid credentials. High severity.

Op-04 is the X.509 certificate chain validation at the peer, which
establishes the authenticity of the client's identity certificate through
the full hierarchy from leaf to trusted root CA. The peer's MSP performs
this validation by checking each certificate signature against its
issuing CA certificate, terminating at a trusted root certificate held in
the peer's local MSP directory. A forged certificate at any level of the
hierarchy enables identity impersonation across the entire channel: any
operation that the forged identity is authorised to perform is available
to the adversary. Op-04 is rated Critical severity because the impact
extends to all channel participants and to all operations governed by
identity-based access controls.

Op-05 is the peer's endorsement signing operation, in which the peer
signs the read-write set produced by chaincode execution with its ECDSA
P-256 private key. The endorsement signature binds the peer's identity
to the specific proposal hash and execution result. It is this signature
that enables the ordering service and committing peers to verify that
a sufficient set of endorsing peers agreed on the transaction's read-write
set. Forgery of this signature enables an adversary to fabricate
endorsements that satisfy any endorsement policy for any proposed
transaction without requiring access to the peer's infrastructure or
executing the chaincode on that peer. High severity.

Op-06 covers all TLS sessions used by the peer for both inbound
connections from clients and SDK applications and outbound connections
to the ordering service and other peers. The peer's long-term TLS
certificate public key is distributed as part of the peer's MSP
configuration and is therefore available as harvest material in the
channel configuration blocks. Medium severity.

\subsection{Phase 3: Transaction Assembly and Submission to Orderer}

Op-07 is the client envelope signing operation, in which the client
re-signs the complete assembled transaction envelope, including the
collected endorsement responses and their signatures, before submitting
it to the ordering service. This re-signing is the final client
authentication check before the transaction enters the consensus
pipeline. Unlike Op-01, which signs only the proposal, Op-07 signs the
complete endorsed transaction, providing a binding that ties the client
identity to the specific set of endorsements collected for the proposal.
Forgery of Op-07 allows the submission of transactions with fabricated
endorsement collections, potentially enabling policy bypass if the
adversary can also forge the endorsement signatures in Op-05. High severity.

Op-08 is the client-to-orderer TLS session established during transaction
submission. The orderer is the network's sequencing and consensus service;
TLS sessions carrying transactions to the orderer include the full
transaction payload, the endorsement signatures, and any client-observable
metadata. Decryption of Op-08 sessions reveals pending transactions before
they are committed to the ledger, which is particularly sensitive in
financial applications where knowledge of pending transactions constitutes
material non-public information. Medium severity.

\subsection{Phase 4: Ordering Service and Block Signing}

Op-09 is the transaction envelope signature verification at the orderer,
which is the final authentication gate before a transaction is sequenced
into a block. High severity.

Op-10 is the block header signing operation by the orderer, in which the
orderer signs the completed block with its ECDSA P-256 private key. This
is the most consequential cryptographic operation in the Fabric
lifecycle. The orderer's block signature is the authoritative assertion
of the canonical transaction sequence; it is what confers on a block its
status as a definitive record of committed state. Forgery of this
signature enables a quantum adversary to inject fraudulent blocks
containing arbitrary transaction data into every peer's ledger. Because
the ledger is append-only and immutable, fraudulent blocks committed
under a forged orderer signature cannot be reversed without out-of-band
channel reconfiguration. Op-10 is rated Critical severity.

Op-11 is the orderer-to-peer TLS session protecting block broadcast.
Medium severity. Op-12 is the mutual TLS authentication between Raft
orderer nodes during the consensus protocol, using ECDSA P-256 for peer
authentication between orderers. Compromise enables man-in-the-middle
attacks on the Raft consensus process. High severity. Op-13 is the Raft inter-orderer TLS session key exchange. It is rated
High severity rather than the Medium typically assigned to ECDH key
exchange operations elsewhere in the inventory, because the channel
this exchange protects carries the Raft consensus protocol traffic
itself: leader-election votes, AppendEntries log replication, and the
ordered set of pending transactions before block commitment.
Decryption of recorded Raft inter-orderer sessions therefore yields
not merely transaction-payload confidentiality loss (as for client- or
peer-facing TLS) but visibility into the consensus state machine
prior to block finalisation, and combined with Op-12 compromise can
enable consensus-protocol manipulation. The integrity-adjacent
character of the leaked information lifts the severity above the
pure-confidentiality threshold used for other ECDH operations.

\subsection{Phase 5: Committing Peer Validation}

Op-14 is the re-verification of every endorsement signature for every
transaction in every received block, performed at every committing peer.
This is the most computationally intensive signature verification
operation in the lifecycle and the operation most directly affected by
the performance implications of post-quantum signature sizes. At a
network throughput of 275 transactions per second, with approximately
two endorsers per transaction and a block batching factor of fifty to
one hundred transactions, this operation processes several thousand
signature verifications per second per peer. The security implications
of Op-14 are High: a fraudulent endorsement signature that is not
detected at Op-14 commits invalid transaction results to the ledger
and invalidates the integrity of the affected ledger state. Op-14 is
rated High rather than Critical because the severity rubric of
Section~\ref{sec:tm_severity} calibrates Critical against the
\textit{scope} of compromise (entire-network integrity or
authenticity) rather than against the \textit{reversibility} of the
consequence. Op-14 enables forgery of endorsement signatures one
identity at a time, with the affected transactions further constrained
by the channel's endorsement policy: an adversary holding a single
recovered endorser key cannot satisfy a multi-organisation policy
without recovering an additional key from a different organisation.
The irreversibility property does apply to the resulting ledger
state, but the same irreversibility applies to every committed
operation in the lifecycle and is therefore not on its own a
distinguishing criterion under the rubric. By contrast, Op-10 (the
orderer block header signing operation) is rated Critical because
compromise of the orderer's signing key permits forgery of the
canonical transaction sequence itself, with a blast radius equal to
every transaction in every block produced under the recovered key, and
without dependence on any per-transaction policy.

The performance dimensions of Op-14 are addressed in detail in
Section~\ref{sec:tm_scale} and are central to the algorithm selection
analysis of Chapter~\ref{ch:empirical}. The critical security
point at Op-14 is that every endorsement signature committed to the
ledger becomes a harvest target. Every endorser identity that has ever
participated in a transaction on the ledger has had its ECDSA P-256
public key embedded in the transaction metadata since the time of the
transaction. An adversary with access to historical ledger data can
recover the endorsement signing key of any peer that has ever endorsed
a transaction, enabling retrospective forgery of endorsements that
purport to have been produced by those peers at any point in the ledger
history. High severity.

Op-15 is the verification of the orderer's block header signature by
each committing peer. This is the final integrity check before the peer
commits the block to its local ledger. Failure to verify this signature
at Op-15 would admit fraudulent blocks into the ledger; correct
verification ensures that only blocks carrying valid orderer signatures
are committed. The severity mirrors Op-10: compromise of the orderer
signing key enables block integrity attacks. Critical severity.

Op-16 is the full certificate chain re-validation performed for every
identity that appears in every transaction in the block, conducted
against the MSP's local trusted root certificates. This operation
ensures that no transaction in the block was signed by an identity
whose certificate has been revoked or whose certificate chain cannot
be validated against a trusted CA. Op-16 represents the committing
peer's final identity authentication check and is the last opportunity
to reject transactions signed by fraudulent certificates before they
are permanently committed. Critical severity.

\subsection{Fabric CA Certificate Infrastructure}

Op-17 is the root CA certificate self-signing operation. The root CA
certificate is the trust anchor for the entire Fabric network. Compromise
of the root CA private key enables a quantum adversary to issue arbitrarily
many valid X.509 certificates for any identity attributes, including
administrator roles, peer node identities, orderer identities, and any
custom attributes used by chaincode access control logic. The root CA
certificate is embedded in the channel genesis block as part of the channel
creation configuration and is replicated to the local MSP directory of
every peer node, making its ECDSA P-256 public key permanently accessible
to any entity that can read the ledger or access a peer's configuration.
Op-17 is rated Critical severity and is the operation of highest strategic
importance in the entire audit. Its strategic importance arises not merely
from the severity of compromise but from the accessibility of the key
material: unlike an orderer block signing key whose public key appears
only in block headers, the root CA public key appears in the very first
block of the ledger and in the MSP configuration of every peer, making
it available for harvest without any ledger scanning.

Op-18 is the intermediate CA certificate signing operation. Most
production Fabric deployments use a two-tier CA hierarchy in which the
root CA signs one or more intermediate CA certificates, and the intermediate
CAs perform the operational certificate issuance to end entities. This
architecture is standard practice in PKI deployments because it reduces
the operational exposure of the root CA key; the root CA can be kept
offline while the intermediate CA handles day-to-day enrolment. From a
quantum threat perspective, however, the intermediate CA certificate
is signed with ECDSA P-256 and its public key appears in the MSP
configuration. An adversary who recovers the root CA private key via
Op-17 can sign fraudulent intermediate CA certificates, and fraudulent
intermediate CA certificates can sign fraudulent end-entity certificates.
The compromise therefore propagates from the root through the full
hierarchy. Op-18 is rated Critical severity.

Op-19 is end-entity certificate issuance to peers, orderers, clients, and
administrators. Each issued certificate contains the enrollee's ECDSA P-256
public key and is signed by the issuing CA with ECDSA P-256. Compromise
of Op-19 is scoped to the entities enrolled under the compromised CA;
it does not automatically propagate to the full network hierarchy in the
way that Op-17 and Op-18 compromise does. High severity. Op-20 is the
Fabric CA TLS session protecting enrolment traffic, including the
transmission of certificate signing requests and issued certificates.
Decryption of Op-20 sessions reveals the identity attributes and
credentials of newly enrolled participants. Medium severity.

\subsection{MSP Identity Validation}

Op-21 is the certificate chain validation performed by the local MSP at
each peer for every identity appearing in the channel's transaction
history. The MSP holds each organisation's trusted root certificate and
validates identity certificate chains against it. This validation is
performed not only at transaction submission time but also at block commit
time for all transactions in the block, making Op-21 a critical continuous
security control. Quantum compromise of any certificate in the hierarchy
enables identity impersonation across the entire organisation: a peer
whose MSP trusts a compromised root CA will accept any identity that
carries a certificate signed under that CA, regardless of whether that
identity was ever legitimately enrolled. Critical severity.

The MSP also maintains a certificate revocation list for each organisation,
checked as part of Op-21. From a quantum threat perspective, certificate
revocation lists are ineffective against TS1 because an adversary with the
root CA private key can issue new, unrevoked certificates at will. Revocation
addresses certificate compromise in the classical threat model, where a
private key is disclosed through theft or operational error; it has no
effect against an adversary who can generate new valid certificates on demand.

Op-22 is the signing of MSP channel configuration updates, which include
channel creation, organisation addition, endorsement policy modification,
and orderer configuration changes. These operations require administrator
ECDSA signatures from one or more organisations, depending on the channel
modification policies. Forgery of administrator signatures at Op-22 enables
arbitrary channel reconfiguration. The most severe consequences include
adding fraudulent organisations with their own CAs to the channel, modifying
endorsement policies to enable single-party fraud, and demoting legitimate
administrators. Critical severity.

\subsection{Gossip Protocol and Key Management}

Op-23 is gossip message signing. Forgery can cause peers to synchronise
to incorrect world state or request unnecessary state transfers. Medium
severity. Op-24 is gossip TLS session key exchange. Decryption of gossip
TLS sessions provides advance knowledge of uncommitted block content,
enabling front-running attacks. Medium severity.

Op-25 is ECDSA P-256 keypair generation by the BCCSP keystore. All
Fabric private keys are generated by this operation. Migration to
post-quantum cryptography requires regeneration of all identity material;
no existing ECDSA key can be converted to a post-quantum equivalent.
High severity.


\subsection{Signature Verification at Scale and the Performance Amplification Effect}
\label{sec:tm_scale}

A crucial observation from the cryptographic inventory concerns the
compounding performance implications of post-quantum signature integration
at Op-14, the endorsement signature re-verification operation. This operation
is performed by every committing peer for every endorsement signature on every
transaction in every received block. In a production network processing 275
transactions per second, with an average of two endorsing peers per
transaction and a block batching factor of fifty to one hundred transactions,
Op-14 processes several thousand signature verifications per second per peer.

The performance cost of post-quantum signatures at Op-14 is amplified at
multiple levels. First, larger post-quantum signatures increase the size of
each transaction record on the ledger, because endorsement signatures are
embedded in the transaction payload. Second, larger transaction records
increase block sizes, because blocks contain the full transaction payloads
of all batched transactions. Third, larger blocks take longer to propagate
across the network via gossip, increasing the interval between the orderer
committing a block and all peers receiving it. Fourth, longer propagation
intervals reduce the sustainable throughput ceiling of the network, because
the rate at which blocks can be safely committed is bounded by the rate at
which they can be delivered to peers for validation.

This chain of effects means that the throughput cost of post-quantum
signature integration is not simply the additional per-signature verification
time, as it would be in a standalone signing system. It is the accumulated
effect of larger signatures, larger blocks, slower propagation, and
constrained throughput, all originating from Op-14's central role in the
commit path. This analysis motivates the expectation, confirmed empirically
in Chapter~\ref{ch:empirical}, that signature output size is a more
important performance determinant than signing speed for blockchain
post-quantum migration. An algorithm that produces smaller signatures will
outperform a faster-signing algorithm with larger signatures in the Fabric
context, even if the faster-signing algorithm would be preferred in a
conventional signing context.

The scale of Op-14 also distinguishes Fabric's performance sensitivity
to post-quantum signatures from that of blockchains with simpler transaction
structures. Each Fabric transaction carries multiple endorsement signatures,
one per required endorser under the endorsement policy, and each endorsement
signature is accompanied by the endorser's full certificate chain. The total
cryptographic payload per transaction can therefore be several times larger
than in a system with a single signature per transaction. This multiplication
factor means that the transition from 72-byte ECDSA signatures to signatures
measured in hundreds or thousands of bytes has a correspondingly amplified
effect on ledger storage and block propagation characteristics.

This relationship can be expressed formally. Let $T_b$ denote the number of
transactions per block, $\bar{e}$ the average number of endorsers per
transaction under the active endorsement policy, $|\sigma_{\mathrm{alg}}|$
the signature size in bytes for algorithm~$\mathrm{alg}$, and
$|\mathit{cert}_{\mathrm{alg}}|$ the certificate chain size. The total
cryptographic payload per block at Op-14 is

\begin{equation}
B_{\sigma} = T_b \cdot \bar{e} \cdot
\bigl(|\sigma_{\mathrm{alg}}| + |\mathit{cert}_{\mathrm{alg}}|\bigr).
\label{eq:block_payload}
\end{equation}

Let $\tau(|\sigma|)$ denote the per-signature verification time, and
$C_{\mathrm{val}}$ the total validation capacity of a committing peer in
operations per second. The maximum sustainable network throughput
$\Gamma_{\max}$ is bounded by

\begin{equation}
\Gamma_{\max} \;\leq\; \frac{C_{\mathrm{val}}}{T_b \cdot \bar{e} \cdot
\tau(|\sigma_{\mathrm{alg}}|)}.
\label{eq:tps_bound}
\end{equation}

Equation~(\ref{eq:tps_bound}) clarifies why $|\sigma_{\mathrm{alg}}|$
is the dominant performance determinant rather than
$\tau(|\sigma_{\mathrm{alg}}|)$ alone. An algorithm with smaller
$|\sigma_{\mathrm{alg}}|$ reduces $B_{\sigma}$ in~(\ref{eq:block_payload}),
decreasing block propagation latency and relaxing the throughput bound
in~(\ref{eq:tps_bound}), even if its individual $\tau(|\sigma|)$ is
higher. This motivates FN-DSA-512 as the preferred algorithm for
throughput-sensitive deployments despite its slower per-operation
signing time relative to ML-DSA-44.

\subsection{Complete Inventory Table}

Figure~\ref{fig:crypto_inventory} maps all twenty-five
quantum-vulnerable operations across the eight component groups of
Hyperledger Fabric v2.5.12. Each operation is displayed as a colour-coded
chip: red for Critical severity, amber for High, and blue for Medium. The
border style of each chip encodes the algorithm family: solid borders
indicate ECDSA digital signature operations and dashed borders indicate
ECDH key exchange operations subject to the harvest-now-decrypt-later
threat through TLS session recording.

Reading the figure from top to bottom traces the complete transaction
lifecycle and then descends into the supporting identity infrastructure.
Phase~1 contributes two operations, one signing and one TLS. Phase~2
is the most cryptographically dense single phase, with four operations
spanning Critical certificate chain validation through Medium TLS session
management. Phase~4 is the most consequential phase for both block
integrity and performance: Op-10, the orderer block header signing
operation, is Critical because its compromise enables fraudulent blocks
to be committed to every peer's immutable ledger, and Op-14, the
endorsement signature re-verification at the committing peer, drives the
performance analysis of Chapter~\ref{ch:empirical} because its throughput
sensitivity to post-quantum signature size is the primary determinant of
network-level performance.

The identity infrastructure rows at the base of the lifecycle section
show the concentration of Critical operations. Four of the six chips in
the Fabric CA and MSP rows are red, confirming that the trust hierarchy
is the most vulnerable architectural component. Notably, Op-17, the root
CA self-signing operation, is the single operation whose compromise
enables complete, network-wide identity forgery through the
harvest-now-decrypt-later attack described in threat scenario TS1 of
Section~\ref{sec:tm_threats}: its public key is embedded in the channel
genesis block and permanently accessible to any ledger reader, making
it the highest-priority migration target.

The dashed-border chips visible across all rows, including Op-02, Op-06,
Op-08, Op-11, Op-13, Op-20, and Op-24, represent ECDH key exchange
operations whose quantum vulnerability manifests through TLS session
decryption rather than signature forgery. Their presence in every phase
confirms that the HNDL threat is not confined to the identity layer but
applies to every inter-component channel in the network.

\begin{figure*}[ht]
\centering
\includesvg[inkscapelatex=false,width=\textwidth]{figures/fig_crypto_inventory}
\caption{Cryptographic inventory of Hyperledger Fabric v2.5.12, showing all twenty-five quantum-vulnerable operations.}
\label{fig:crypto_inventory}
\end{figure*}

Table~\ref{tab:crypto_inventory} presents the complete inventory of all
twenty-five quantum-vulnerable operations identified in this audit. The
table confirms that the vulnerability is both complete, covering every
public-key cryptographic operation in the standard lifecycle, and
uniform, with every operation relying on ECDLP-based mathematics. The
distribution across severity levels shows eight Critical operations,
ten High operations, and seven Medium operations. The Critical operations
are concentrated in the certificate infrastructure, the block signing
and verification path, and the MSP validation layer, confirming that the
integrity backbone of the network is most directly threatened.

Figure~\ref{fig:threat_model_overview} provides a component-level
overview of the vulnerability landscape. The diagram organises the six
principal Fabric components into two horizontal tiers. The upper tier
shows the four transaction lifecycle components, progressing left to
right through the phases of proposal, endorsement, ordering, and
commitment. The lower tier shows the identity management infrastructure,
comprising Fabric CA and the MSP, which serve all four lifecycle
components, together with the gossip protocol layer.

Color encoding maps directly to the severity classification established
in Section~\ref{sec:tm_method}. Red components hold at least one
Critical-severity operation whose compromise enables network-wide
integrity failure. The endorsing peer, ordering service, and committing
peer are all red because each contains block signing or certificate
chain validation operations rated Critical. Fabric CA and the MSP are
also red, and their grouping within the dashed red overlay identifies
the TS1 attack surface: a single quantum key recovery against the root
CA's ECDSA P-256 public key, which is embedded in the genesis block and
permanently accessible to any channel participant, is sufficient to
compromise every identity in the network. The client SDK is amber,
reflecting High-severity operations whose impact is bounded to
transactions that the compromised identity is authorised to submit.
The gossip protocol is blue, reflecting Medium-severity confidentiality
threats that require network observation rather than passive ledger
access.

The dashed blue line running the full width of the diagram represents
the TLS session layer corresponding to threat scenario TS4. Every
connection between any pair of the six components uses TLS 1.3 with
ECDH P-256 key exchange, and every recorded TLS handshake is
retrospectively decryptable by a quantum adversary once a
cryptographically relevant quantum computer becomes available. The
threat is classified Medium rather than Critical because it targets
confidentiality rather than ledger integrity, and because exploitation
requires prior recording of network traffic rather than mere read
access to the public ledger.

The five threat scenario summary cards at the base of the diagram
correspond to the five formal threat scenarios developed in
Section~\ref{sec:tm_threats}. Their left-to-right arrangement
reflects the risk prioritisation of Section~\ref{sec:tm_risk}: TS1
and TS3 at Critical appear first, TS2 at High follows, and TS4 and
TS5 at Medium appear last. The TS2 badge placed directly on the
endorser-to-orderer flow arrow indicates that transaction signature
forgery is possible wherever transaction signatures and their
corresponding public keys are committed to the ledger, which is at
every lifecycle phase.

\begin{figure*}[ht]
\centering
\includesvg[inkscapelatex=false,width=\textwidth]{figures/fig_threat_model}
\caption{Quantum vulnerability overview for Hyperledger Fabric v2.5.12, colour-coded by severity across six components.}
\label{fig:threat_model_overview}
\end{figure*}


Table~\ref{tab:crypto_inventory} lists all twenty-five operations
individually, specifying the precise algorithm and key type for each
and recording the severity assessment. The table is organised by
the same component groupings shown in
Figure~\ref{fig:threat_model_overview}, allowing the reader to trace
each severity assessment in the diagram back to the specific operations
that produce it.

\begin{table*}[ht]
\centering
\caption{Complete cryptographic inventory of Hyperledger Fabric
v2.5.12. All twenty-five operations use ECDLP-based mathematics and
are threatened by Shor's algorithm. Severity codes: C~=~Critical,
H~=~High, M~=~Medium.}
\label{tab:crypto_inventory}
\begin{tabular}{@{}p{0.75cm}>{\raggedright\arraybackslash}p{5.3cm}%
  >{\raggedright\arraybackslash}p{3.8cm}c@{}}
\toprule
\textbf{Op} & \textbf{Operation} & \textbf{Algorithm} &
\textbf{Sev.} \\
\midrule
\multicolumn{4}{@{}l}{\textbf{Phase 1 : Transaction Proposal}} \\[2pt]
01 & Client transaction proposal signing   & ECDSA P-256 & H \\
02 & Client TLS handshake to peer          & ECDH P-256  & M \\[5pt]
\multicolumn{4}{@{}l}{\textbf{Phase 2 : Peer Endorsement}} \\[2pt]
03 & Client signature verification at peer & ECDSA P-256 & H \\
04 & X.509 certificate chain validation    & ECDSA P-256 & C \\
05 & Chaincode execution result signing    & ECDSA P-256 & H \\
06 & Peer TLS inbound/outbound sessions    & ECDH P-256  & M \\[5pt]
\multicolumn{4}{@{}l}{\textbf{Phase 3 : Transaction Assembly}} \\[2pt]
07 & Client transaction envelope signing   & ECDSA P-256 & H \\
08 & Client-to-orderer TLS session         & ECDH P-256  & M \\[5pt]
\multicolumn{4}{@{}l}{\textbf{Phase 4 : Ordering Service}} \\[2pt]
09 & Envelope signature verification at orderer & ECDSA P-256 & H \\
10 & Block header signing by orderer       & ECDSA P-256 & C \\
11 & Orderer-to-peer TLS (block broadcast) & ECDH P-256  & M \\
12 & Raft inter-orderer mutual TLS auth    & ECDSA P-256 & H \\
13 & Raft inter-orderer TLS key exchange   & ECDH P-256  & H \\[5pt]
\multicolumn{4}{@{}l}{\textbf{Phase 5 : Committing Peer Validation}} \\[2pt]
14 & Endorsement signature re-verification & ECDSA P-256 & H \\
15 & Block header signature verification   & ECDSA P-256 & C \\
16 & Certificate chain re-validation       & ECDSA P-256 & C \\[5pt]
\multicolumn{4}{@{}l}{\textbf{Fabric CA: Certificate Infrastructure}} \\[2pt]
17 & Root CA certificate self-signing      & ECDSA P-256 & C \\
18 & Intermediate CA certificate signing   & ECDSA P-256 & C \\
19 & End-entity certificate issuance       & ECDSA P-256 & H \\
20 & Fabric CA TLS session                 & ECDH P-256  & M \\[5pt]
\multicolumn{4}{@{}l}{\textbf{MSP: Identity Validation}} \\[2pt]
21 & MSP certificate chain validation      & ECDSA P-256 & C \\
22 & MSP channel configuration signing     & ECDSA P-256 & C \\[5pt]
\multicolumn{4}{@{}l}{\textbf{Gossip Protocol}} \\[2pt]
23 & Gossip message signing                & ECDSA P-256 & M \\
24 & Gossip TLS sessions                   & ECDH P-256  & M \\[5pt]
\multicolumn{4}{@{}l}{\textbf{Key Management}} \\[2pt]
25 & BCCSP keypair generation              & ECDSA P-256 & H \\
\bottomrule
\end{tabular}
\end{table*}


\section{Quantum Threat Model}
\label{sec:tm_threats}

\subsection{Adversarial Capability Model}

The threat model assumes an adversary who possesses a cryptographically
relevant quantum computer: a fault-tolerant quantum processor capable of
executing Shor's algorithm against P-256 at the key sizes used by Fabric.
Roetteler et al.\ established that breaking ECDSA P-256 requires
approximately 2,330 logical qubits with appropriate error correction
overhead~\cite{roetteler2017}. Gidney and Eker{\aa} provided analogous
estimates for RSA-2048~\cite{gidney2021}. While no such device currently
exists, the harvest-now-decrypt-later threat requires it to be treated as
operationally relevant today for data whose required security lifetime
extends beyond the projected CRQC emergence window~\cite{mosca2018}.

The adversary is assumed to possess the following classical capabilities
in addition to the quantum processor: passive observation of Fabric
network traffic including TLS handshakes and gossip messages; read access
to the public ledger including all committed blocks and channel
configuration blocks; and the ability to query the Fabric ledger for
historical transaction data. These capabilities require no privileged
access and are available to any channel participant. The adversary is
assumed not to have physical access to Fabric nodes, prior knowledge of
private keys, or the ability to corrupt a majority of Raft orderer nodes
through classical means.

\subsection{TS1: Identity Impersonation via Certificate Forgery}

Severity: Critical. Exploitability: High. Priority: P1.

TS1 targets the root CA certificate embedded in the channel genesis block.
The pre-conditions for this attack are the ability to read any channel
genesis block or MSP configuration, which is available to all channel
participants, and possession of a CRQC capable of Shor's algorithm.
Neither condition requires privileged access.

The adversary reads the genesis block or channel configuration to obtain
the root CA X.509 certificate. The certificate's SubjectPublicKeyInfo
field contains the CA's ECDSA P-256 public key. The adversary applies
Shor's algorithm to recover the root CA private key. With this key,
the adversary issues arbitrarily many valid certificates for any identity
attributes, including administrator roles, peer node identities, orderer
identities, and elevated client identities.

The impact is total. Every certificate issued by the compromised CA, and
every transaction signed by an identity holding such a certificate, passes
MSP validation and is accepted by all honest peers. The adversary can also
sign fraudulent channel configuration updates with forged administrator
credentials, enabling modification of endorsement policies, addition of
fraudulent organisations, and revocation of legitimate participants.

The harvest-now-decrypt-later dimension is permanent. The root CA
certificate is embedded in the genesis block and replicated to all peer
MSP directories. A channel operational since 2018 has had its root CA
public key continuously available as harvest material for years. Post-quantum
migration of future certificates does not remove the historical CA public
key from the ledger.

Affected operations: Op-04, Op-16, Op-17, Op-18, Op-19, Op-21, Op-22.

The severity of TS1 is further compounded by the chain-of-trust structure
of Fabric's certificate hierarchy. In most production deployments, the
root CA signs one or more intermediate CA certificates, and the intermediate
CAs in turn sign end-entity certificates for peers, orderers, clients, and
administrators. An adversary who recovers the root CA private key can sign
fraudulent intermediate CA certificates, and fraudulent intermediate CA
certificates can in turn sign fraudulent end-entity certificates for any
identity. The full certificate hierarchy is compromised from a single
point of quantum attack. This attack does not require the adversary to
compromise each individual peer or orderer key; the trust hierarchy
propagates the compromise from a single root key recovery to all
identities in the network.

Furthermore, the MSP channel configuration signing operation Op-22 is
within the scope of TS1 because MSP configuration updates are signed by
administrator identities whose certificates chain to the CA. An adversary
holding fraudulent administrator certificates can submit channel
configuration updates that modify endorsement policies, add fraudulent
organisations with their own CAs, or revoke the certificates of legitimate
participants. These channel configuration attacks are particularly
dangerous because they are cryptographically indistinguishable from
legitimate administrative actions: they carry valid signatures from
valid-looking certificates signed by what appears to be a valid CA.

\subsection{TS2: Transaction Signature Forgery}

Severity: High. Exploitability: High. Priority: P2.

TS2 targets ECDSA P-256 signatures on transaction proposals, endorsement
responses, and transaction envelopes. Transaction signatures and the public
keys required to verify them are embedded in every transaction record on
the Fabric ledger. Any entity with ledger read access can extract any
participant's ECDSA public key from historical transactions. The adversary
applies Shor's algorithm to a harvested public key to recover the
corresponding private key and signs arbitrary transactions on behalf of
the compromised identity.

The impact depends on the compromised identity. A client identity
compromise is scoped to transactions that identity is authorised to
submit. An endorsing peer identity compromise extends to fabricating
endorsements satisfying the endorsement policy of any chaincode deployed
on that peer, enabling policy bypass without access to peer infrastructure.
An orderer identity compromise encompasses block signing as described in TS3.

Every historical transaction on the Fabric ledger contains the ECDSA
public key of its signatories as part of the transaction metadata. Any
entity with ledger read access has implicitly harvested these keys. A
quantum adversary who acquires any portion of the ledger history can
recover the private key of any identity that has signed a transaction on
that ledger.

The harvest-now-decrypt-later dimension of TS2 is permanent in the same
way as TS1. The ECDSA public keys of all historical transaction
signatories are immutably embedded in the ledger. A network that migrated
to post-quantum signing today would protect all future transaction
signatures but cannot remove the historical ECDSA public keys of
participants who signed transactions in the past. This means that the
private keys of all historical transaction signatories remain
retrospectively recoverable for any adversary who holds a copy of the
ledger and gains access to a CRQC.

The operational impact of TS2 is further compounded by the role of
endorsement policies in Fabric's security model. An endorsement policy
specifies which set of organisations or identities must endorse a
transaction for it to be considered valid. An adversary who recovers
the signing keys of identities across different organisations can
fabricate endorsements that satisfy multi-organisation endorsement
policies, enabling fraudulent transactions to pass the endorsement check
even for chaincodes with strict cross-organisational approval requirements.
For supply chain networks, financial settlement networks, or any
application where the endorsement policy implements a business rule
requiring independent confirmation from multiple parties, TS2 enables
complete bypass of the multi-party approval mechanism.

Affected operations: Op-01, Op-03, Op-05, Op-07, Op-09, Op-14.

\subsection{TS3: Block Integrity Compromise via Orderer Key Recovery}

Severity: Critical. Exploitability: High. Priority: P1.

TS3 targets the ECDSA P-256 block signing keys of orderer nodes. The
orderer's block signatures and corresponding public keys appear in the
header of every block on the Fabric ledger, accessible to all channel
participants. The adversary extracts an orderer public key from any block
header and applies Shor's algorithm to recover the private key. With this
key, the adversary produces valid block signatures for fraudulent blocks
of arbitrary content.

The impact is Critical. The orderer's block signature is the authoritative
assertion of the canonical transaction sequence. A fraudulent block
carrying a valid orderer signature is accepted by all committing peers as
a legitimate ledger record, enabling injection of false transactions,
reordering of transactions, or suppression of legitimate activity.
Because the ledger is append-only and immutable, fraudulent blocks
committed under a recovered orderer signature cannot be reversed without
out-of-band channel reconfiguration. This threat is of particular concern
in financial settlement applications where transaction ordering determines
settlement outcomes.

Affected operations: Op-10, Op-12, Op-13, Op-15.

The irreversibility of committed ledger records is the defining
characteristic that elevates TS3 from a High severity to a Critical severity
threat. In a conventional database system, an unauthorised record can be
identified and corrected through administrative procedures. In a Fabric
blockchain, a block that has been committed to the ledger by a committing
peer and acknowledged by the gossip dissemination protocol is permanent.
It forms part of the hash-chained sequence that subsequent blocks reference
in their previous-hash field. Removing or modifying a committed block would
invalidate the hash chain for all subsequent blocks and require
reconstruction of the entire ledger history, an operation that is not
supported by any standard Fabric administrative procedure.

This irreversibility means that the damage from TS3 is not merely
temporary. A fraudulent block committed under a forged orderer signature
becomes a permanent part of every peer's ledger. Even if the attack is
detected and the network is subsequently reconfigured with new orderer
keys, the fraudulent block remains in the historical ledger. Organisations
relying on the blockchain as an audit trail, a settlement record, or a
regulatory compliance archive will find the integrity of that archive
permanently compromised for the period during which the fraudulent blocks
were committed.

The Raft inter-orderer operations Op-12 and Op-13 add a further dimension
to TS3. The Raft consensus protocol requires a majority of orderer nodes
to agree on the content of each log entry before it is committed as a
block. An adversary who compromises Raft inter-orderer authentication
can perform man-in-the-middle attacks on the consensus protocol itself,
potentially influencing which transactions are included in which blocks
or disrupting the leader election process. While this attack is more
complex than simple block header forgery, it could enable selective
transaction suppression or ordering manipulation that is difficult to
detect even with correct block signatures from legitimate orderers.

\subsection{TS4: TLS Session Decryption via Harvest-Now-Decrypt-Later}

Severity: Medium. Exploitability: Medium. Priority: P3.

TS4 targets ECDH P-256 key exchange in all TLS sessions between Fabric
components. An adversary positioned to observe Fabric network traffic
records TLS handshakes and encrypted session content. The ECDH key
exchange parameters appear in cleartext in the TLS handshake. When a
CRQC becomes available, the adversary applies Shor's algorithm to the
recorded ECDH exchange to recover the session key and decrypts all
previously recorded traffic.

The impact is primarily a confidentiality breach. Decrypted sessions
reveal transaction proposals, endorsement responses, block delivery
content, and administrative communications prior to their public ledger
commitment. For channels carrying commercially sensitive or personally
identifiable data, this represents significant regulatory and commercial
exposure.

An important technical distinction applies to forward secrecy. TLS
deployments commonly use ephemeral ECDH to provide forward secrecy
against classical adversaries. Against a quantum adversary, ephemeral
ECDH provides no forward secrecy, because the ephemeral exchange itself
is quantum-vulnerable. A quantum adversary who records the TLS handshake
recovers the ephemeral session key directly without needing the long-term
certificate private key. TLS does not provide forward secrecy against a
quantum adversary in this model.

Affected operations: Op-02, Op-06, Op-08, Op-11, Op-13, Op-20, Op-24.

The data categories most affected by TS4 vary by deployment context.
In a healthcare Fabric network, TLS sessions carry transaction proposals
that include patient identifiers, diagnosis codes, prescription data, and
treatment records before these are committed to the ledger. A quantum
adversary who decrypts these sessions has access to the same sensitive
health information as a party who can read the committed ledger records,
but potentially earlier and without the access controls that govern ledger
read permissions. In a financial settlement network, TLS sessions carry
transaction proposals whose details, including counterparties, amounts,
and settlement timing, may be market-sensitive information that should
not be available to third parties prior to settlement finalisation.

The Fabric CA enrolment TLS sessions covered by Op-20 are also within
TS4 scope. Enrolment sessions carry the credentials and identity
attributes of newly enrolling participants, including administrator
accounts, new peer node identities, and client application credentials.
Decryption of enrolment sessions provides an adversary with a catalogue
of all identities in the network, their roles, and the timing of their
addition, information that is relevant to understanding the network's
operational history and governance structure.

\subsection{TS5: Gossip Protocol Manipulation}

Severity: Medium. Exploitability: Low. Priority: P4.

TS5 targets ECDSA gossip message signatures and ECDH gossip TLS sessions.
The adversary recovers a peer's gossip signing key from observable gossip
protocol messages, in which signing keys and public key certificates are
distributed as part of peer discovery, and injects fraudulent gossip
messages signed with the recovered key. Alternatively, the adversary
records gossip TLS sessions and decrypts them when a CRQC becomes
available.

Two distinct impact categories arise from TS5. The first is state
synchronisation manipulation. Gossip message forgery can cause peers to
accept false announcements about the availability of new blocks, request
unnecessary state transfers, accept incorrect ledger height information,
or apply an incorrect world state from a state transfer initiated by a
fraudulent gossip message. In large networks where state transfer from
trusted peers is the recovery mechanism for nodes that have been offline,
TS5 could cause a recovering peer to restore from a fraudulent state
rather than from the correct ledger history.

The second impact category is transaction content exposure. Gossip
protocol messages disseminate blocks to peers who have not yet received
them directly from the orderer. An adversary who can decrypt gossip TLS
sessions observes block content as it propagates across the network,
potentially obtaining advance knowledge of transaction results before
the affected parties have been notified through normal application
channels. For financial networks processing high-value transactions,
this advance knowledge constitutes material non-public information with
direct economic value.

The lower exploitability of TS5 relative to TS1 through TS3 reflects
the requirement for network observation or message injection capabilities.
Unlike the P1 and P2 threats, which can be executed by any entity with
read access to the public ledger, TS5 requires the adversary to observe
or inject gossip protocol traffic, which requires a network position
that the adversary may not hold. This additional requirement lowers the
exploitability rating while not eliminating the threat for adversaries
who have previously been channel participants or who have observed
network traffic from other vantage points.

Affected operations: Op-23, Op-24.

Figure~\ref{fig:threat_attack_paths} summarises the five threat scenarios
as a matrix relating each scenario to the Fabric component it targets
(entry point, shown by the downward arrow) and the components it impacts
(filled circles). The two Critical-severity scenarios, TS1 and TS3, both
originate from components whose key material is permanently accessible on
the immutable ledger: the Fabric CA root certificate in the genesis block
(TS1) and the orderer signing key reconstructable from historical block
headers (TS3). TS2 originates at the BCCSP signing layer and affects
transaction records in the ledger. TS4 and TS5 require active network
observation rather than ledger access and are therefore rated Medium
severity.

\begin{figure*}[ht]
\centering
\includesvg[inkscapelatex=false,width=\textwidth]{figures/fig_threat_attack_paths}
\caption[Threat scenario attack paths mapped to Fabric components.]
{Threat scenario attack paths mapped to Fabric components: entry points and impacted components.}
\label{fig:threat_attack_paths}
\end{figure*}

\section{Risk Prioritisation}
\label{sec:tm_risk}

\subsection{Risk Matrix}

Table~\ref{tab:risk_matrix} presents the risk prioritisation matrix,
assessing each threat scenario against two dimensions: severity and
exploitability. The combined assessment determines the migration priority.

\begin{table*}[ht]
\centering
\caption{Risk prioritisation matrix for the five quantum threat scenarios
identified in the Hyperledger Fabric v2.5.12 audit.}
\label{tab:risk_matrix}
\setlength{\tabcolsep}{5pt}
\footnotesize
\begin{tabular}{@{}>{\raggedright\arraybackslash}p{1.6cm}>{\raggedright\arraybackslash}p{5.0cm}cc>{\raggedright\arraybackslash}p{3.5cm}@{}}
\toprule
\textbf{Priority} & \textbf{Scenario} & \textbf{Severity} &
\textbf{Exploit.} & \textbf{Rationale} \\
\midrule
P1 & TS1: Identity Impersonation & Critical & High &
  CA certificate in public genesis block \\
P1 & TS3: Block Integrity Compromise & Critical & High &
  Orderer key in every block header \\
P2 & TS2: Transaction Signature Forgery & High & High &
  Tx signatures in every ledger record \\
P3 & TS4: TLS Session Decryption & Medium & Medium &
  Requires network observation \\
P4 & TS5: Gossip Protocol Manipulation & Medium & Low &
  Requires active network position \\
\bottomrule
\end{tabular}
\end{table*}

Table~\ref{tab:risk_matrix} shows that TS1 and TS3 share the joint
highest priority, both rated Critical severity and High exploitability.
TS2 is rated second priority, with High severity and High exploitability.
TS4 and TS5 occupy the lower priorities, constrained by network
observation requirements that are more demanding than passive ledger
reading.

\subsection{Architectural Conclusions}

Three architectural conclusions follow from the risk matrix. The first
is that the CA certificate infrastructure and the block signing path
share the top priority and must be addressed simultaneously. Both attack
surfaces are accessible to any passive ledger observer because the CA
public key is in the genesis block and the orderer public key is in every
block header. No defensive measure short of replacing both with
post-quantum equivalents can address the P1 threats.

The second conclusion is that peer endorsement and client transaction
signing share second priority. Their impact is more limited in scope than
TS1 and TS3, but exploitability is equally high because transaction
signatures and public keys appear in every historical transaction record.

The third conclusion is that TLS and gossip operations occupy lower
priority, not because their protection is unimportant but because their
exploitation requires network observation access rather than mere ledger
access. These operations are nonetheless included in the integration
requirements because a complete quantum resistance posture requires
protection of the full attack surface, and because HNDL attacks on TLS
require only past network observation, not ongoing access.

\subsection{Quantitative Risk Synthesis}
\label{sec:tm_risk_quant}

The qualitative ratings in Table~\ref{tab:risk_matrix} can be
formalised into a composite risk score. Map severity to
$s_i \in \{1, 2, 3\}$ (Medium~$= 1$, High~$= 2$, Critical~$= 3$)
and exploitability to $e_i \in \{1, 2, 3\}$ (Low~$= 1$,
Medium~$= 2$, High~$= 3$). The composite risk score is then

\begin{equation}
\rho_i = s_i \cdot e_i.
\label{eq:risk_score}
\end{equation}

Evaluating equation~(\ref{eq:risk_score}) for each threat scenario:
TS1 and TS3 both yield $\rho = 3 \times 3 = 9$ (P1, maximum
attainable score); TS2 yields $\rho = 2 \times 3 = 6$ (P2);
TS4 yields $\rho = 1 \times 2 = 2$ (P3); and TS5 yields
$\rho = 1 \times 1 = 1$ (P4). No defensive measure operating
beneath the cryptographic layer can reduce $s_i$ or $e_i$
for TS1 or TS3, because their exploitability derives from
immutable ledger content.

The completeness of the requirements set with respect to the threat
model can be stated formally. Define the coverage relation
$\Phi \subseteq \{R_1,\ldots,R_8\} \times \{\mathrm{TS1},\ldots,\mathrm{TS5}\}$
where $(R_k, \mathrm{TS}_j) \in \Phi$ if and only if requirement $R_k$
addresses threat scenario $\mathrm{TS}_j$. The requirement set provides
complete threat coverage if and only if

\begin{equation}
\forall\, j \in \{1,\ldots,5\},\;\;
\exists\, k \in \{1,\ldots,8\} \;:\; (R_k,\,\mathrm{TS}_j) \in \Phi.
\label{eq:req_coverage}
\end{equation}

Inspection of Table~\ref{tab:req_traceability} confirms that
(\ref{eq:req_coverage}) holds: TS1 is covered by $R_3$, TS2 and TS3
by $R_1$, and TS4 and TS5 by $R_2$. The requirement set is therefore
complete with respect to the threat scenarios identified in
Section~\ref{sec:tm_threats}; this completeness statement is bounded
by the scope of the threat enumeration and does not rule out future
discovery of threats outside the present taxonomy.

Finally, the rate at which new ECDSA key material enters the permanent
harvest set, and hence the rate at which HNDL exposure accumulates,
is a function of current network throughput. Let $\Gamma(t)$ denote
the network throughput in transactions per second at time $t$,
$\bar{k}$ the mean number of distinct new ECDSA public keys committed
per transaction (signing identity keys plus certificate chain keys),
and $\mathbb{1}[t < t_m]$ the indicator of the pre-migration period.
The instantaneous growth rate of the harvest set is

\begin{equation}
\frac{d\,|\mathcal{H}|}{dt}(t)
= \Gamma(t)\cdot\bar{k}\cdot\mathbb{1}[t < t_m].
\label{eq:hndl_rate}
\end{equation}

Integrating equation~(\ref{eq:hndl_rate}) over the network's full
operational history from genesis at $t_0$ to the evaluation date
$t_{\mathrm{eval}}$ gives the total accumulated harvest exposure:

\begin{equation}
|\mathcal{H}(t_{\mathrm{eval}})| =
\int_{t_0}^{\min(t_m,\,t_{\mathrm{eval}})}
\Gamma(t)\,\bar{k}\;\mathrm{d}t
\;\approx\; \bar{\Gamma}\,\bar{k}\,
(t_{\mathrm{eval}} - t_0),
\label{eq:hndl_total}
\end{equation}

where the approximation on the right holds for a network with
approximately constant mean throughput $\bar{\Gamma}$ that has not
yet completed migration. Equation~(\ref{eq:hndl_total}) makes precise
the governance conclusion of Section~\ref{sec:tm_discussion}: every
day of continued operation under ECDSA P-256 permanently enlarges the
harvest set by $\approx \bar{\Gamma}\cdot\bar{k}\cdot 86{,}400$ keys.

\subsection{The Partial Migration Vulnerability}
\label{sec:tm_partial}

A critical insight emerging from the risk matrix is the structural
interdependence between the CA certificate layer and the BCCSP signing
layer. If an operator deploys post-quantum signing at the BCCSP layer
without simultaneously migrating the certificate infrastructure, the CA
public keys in the channel configuration blocks remain in ECDSA form.
A quantum adversary who recovers the CA private key from these ECDSA
public keys can issue fraudulent post-quantum certificates accepted as
legitimate by any PQC-capable MSP. The quantum resistance of the signing
layer is completely neutralised by the classical vulnerability of the
certificate layer.

Figure~\ref{fig:partial_migration} contrasts the two deployment states
directly. The colour coding encodes migration status throughout both panels:
teal components have been migrated to post-quantum algorithms and are secure;
red components retain ECDSA P-256 and are quantum-vulnerable.

In the left panel, the BCCSP signing layer and the MSP are both teal. The
BCCSP layer now signs transactions with ML-DSA or FN-DSA, and the MSP runs
the post-quantum \texttt{buildPQCChain()} certificate chain validator. However,
the Fabric CA box is red: it still issues certificates under ECDSA P-256.
Because the CA is ECDSA, the CA public key committed to the channel genesis
block is also an ECDSA P-256 key. This is shown by the ledger pill at the base
of the left column, labelled ``Ledger: ECDSA P-256 CA key in genesis block.''
That public key is permanent and accessible to every channel participant.

The four small boxes to the right of the left column trace the attack steps
that TS1 executes against this configuration. A dashed arc connects the
ledger pill to the first attack box, representing the adversary reading the
ECDSA CA public key from the immutable ledger. Step~1 applies Shor's algorithm
on a cryptographically relevant quantum computer to that harvested public key.
Step~2 produces the CA private key from the computation. Step~3 uses the
recovered private key to sign a fraudulent post-quantum certificate. The arrow
at the bottom of the attack chain then points back into the MSP box, because
step~4 is that MSP accepting the fraudulent certificate as legitimate: the
certificate carries a valid post-quantum signature produced under the real CA
key, and the PQC chain validator has no basis on which to reject it.

The critical observation is that the teal BCCSP and MSP boxes play no role
in this attack. The adversary never interacts with them. The entire attack path
runs from the ledger to the CA, bypassing both post-quantum layers entirely.
Deploying post-quantum signing at the BCCSP layer provides no protection
because the attack does not target transaction signatures; it targets the CA
key that authorises identities. The post-quantum MSP actively assists the
attack by successfully validating the forged certificate.

The right panel shows the complete \mfabric{} deployment. All three layer
boxes are teal, and the ledger pill is also teal: the Fabric CA now satisfies
requirement~R3 by issuing certificates under an ML-DSA or FN-DSA key, so the
genesis block contains a post-quantum CA public key. The margin annotation
``No ECDSA key to recover'' marks the point where the attack would begin in
the left panel. There is no ECDSA public key to harvest. Shor's algorithm
cannot recover a private key from a lattice-based public key. The attack
has no entry point and the result box reads ``TS1 closed.''

This interdependence is formalised in Section~\ref{sec:tm_requirements}:
threat scenario TS1 is addressed exclusively by the certificate
infrastructure requirement R3 and cannot be addressed by any combination
of BCCSP signing requirements. An operator who completes only the BCCSP
signing migration and defers the certificate infrastructure migration
has an incomplete quantum resistance posture regardless of which
post-quantum signing algorithms are deployed.

\begin{figure*}[ht]
\centering
\includesvg[inkscapelatex=false,width=\textwidth]{figures/fig_partial_migration}
\caption{Partial migration vulnerability: TS1 remains fully exploitable when only the BCCSP layer is migrated.}
\label{fig:partial_migration}
\end{figure*}

\section{Quantum Attack Timeline Analysis}
\label{sec:tm_mosca}

\subsection{Mosca's Inequality}

The urgency of post-quantum migration cannot be established by threat
severity alone; it requires a temporal analysis relating the security
lifetime of protected data to the anticipated CRQC emergence date and the
migration lead time. Mosca's inequality provides this framework~\cite{mosca2018}.
The inequality states that if $X + Y > Z$ then migration must begin
immediately, where $X$ is the number of years before a CRQC becomes
operational, $Y$ is the number of years required to complete migration,
and $Z$ is the number of years for which the protected data must remain
secure. Formally:

\begin{equation}
X + Y > Z \;\Longrightarrow\; \text{post-quantum migration must begin immediately.}
\label{eq:mosca}
\end{equation}

For this analysis, two CRQC timeline estimates are used, anchored to a
common planning base year of 2026 (the submission year of this
thesis). The optimistic estimate places $X$ at 7~years, corresponding
to a CRQC available approximately in 2033 and aligning with the
optimistic marker shown in Figure~\ref{fig:hndl_timeline} of
Chapter~\ref{ch:introduction}. This estimate is informed by Mosca's
2018 analysis, which assessed a one-in-seven chance of a CRQC by 2026
and a one-in-two chance by 2031~\cite{mosca2018}; with those
probability milestones now in or behind the present planning horizon,
the next median-arrival estimate used in this thesis is 2033. The
pessimistic estimate places $X$ at 15~years, corresponding to a CRQC
available approximately in 2041, which represents a more conservative
engineering-based projection. The migration time $Y$ for an enterprise Fabric network is
estimated at 2 to 3 years. This estimate accounts for: the software
development and testing cycle required to produce a production-quality
post-quantum Fabric release; the certificate infrastructure renewal
process, which requires all enrolled identities to obtain new
post-quantum certificates; the staged multi-organisation rollout that
is necessary in permissioned networks where upgrade coordination spans
organisational boundaries; and the operational validation and regulatory
approval processes that govern production deployment changes in regulated
industries. The estimate is consistent with the timelines adopted by
major financial sector infrastructure providers in their post-quantum
migration planning documentation.

\subsection{Application to Fabric Deployment Scenarios}

Table~\ref{tab:mosca} presents the application of Mosca's inequality to
five representative Fabric deployment scenarios. For each scenario, the
required security lifetime $Z$ is estimated and the inequality is
evaluated under both CRQC timeline estimates.

\begin{table}[ht]
\centering
\caption{Application of Mosca's inequality to five representative
Hyperledger Fabric deployment scenarios, evaluated under optimistic
and pessimistic CRQC timeline estimates.}
\label{tab:mosca}
\setlength{\tabcolsep}{5pt}
\footnotesize
\begin{tabular}{@{}
>{\raggedright\arraybackslash}p{2.2cm}
>{\centering\arraybackslash}p{1.2cm}
>{\centering\arraybackslash}p{0.5cm}
>{\centering\arraybackslash}p{0.5cm}
>{\centering\arraybackslash}p{0.5cm}
>{\raggedright\arraybackslash}p{2.0cm}
@{}}
\toprule
\textbf{Scenario} & \textbf{Z~(yr)} &
\textbf{X~opt} & \textbf{X~pess} & \textbf{Y} &
\textbf{Action} \\
\midrule
Healthcare records   & 30+   & 7 & 15 & 3 & Immediate \\
Financial settlement & 20+   & 7 & 15 & 3 & Immediate \\
Govt.\procurement   & 25+   & 7 & 15 & 3 & Immediate \\
Supply chain         & 5--10 & 7 & 15 & 3 & Borderline/Immediate \\
Two-year contracts   & 2     & 7 & 15 & 3 & Can defer \\
\bottomrule
\end{tabular}
\end{table}

Table~\ref{tab:mosca} shows that three of the five scenarios satisfy
Mosca's inequality under both timeline estimates, indicating immediate
migration action regardless of when the CRQC materialises. Healthcare
networks with patient data subject to thirty-year or longer regulatory
retention requirements, financial settlement networks recording
transactions whose integrity must be verifiable for at least twenty
years, and government procurement networks with multi-decade archival
obligations all require immediate migration. The supply chain entry is
labelled \textit{Borderline / Immediate} because supply chain ledger
records span a wide sensitivity range: short-lived inventory records
may carry $Z\!\approx\!5$ years while traceability and provenance
records driven by food-safety, pharmaceutical, or trade-compliance
regulation may carry $Z\!\approx\!10$ years. Under the optimistic
timeline ($X{=}7$, $Y{=}3$, $X{+}Y{=}10$), the inequality $X{+}Y\!>\!Z$
is satisfied for the lower end of the range and fails to be strictly
satisfied at $Z{=}10$; under the pessimistic timeline ($X{+}Y{=}18$)
the inequality is satisfied for the entire range. The action label is
therefore conditional on the operator's view of the sensitivity
lifetime of the records being protected. Only applications with data
sensitivity lifetimes of two years or fewer fall safely below the
$X{+}Y$ horizon under both timelines.

\subsection{HNDL Amplification in Blockchain Deployments}

The conventional application of Mosca's inequality evaluates $Z$ against
the desired future security lifetime. For blockchain deployments, this
is insufficient. The append-only ledger means that every record
previously committed is also a harvest target. The correct blockchain
formulation of the inequality substitutes for $Z$ the maximum sensitivity
lifetime of any record currently on the ledger, which for healthcare,
financial, and government networks is effectively unbounded.

Let $t_0$ denote the network genesis date, $t_m$ the migration completion
date, and $\mathcal{L}(t)$ the set of public keys committed to the ledger
up to time $t$. The HNDL exposure set at any time $t \geq t_0$ is

\begin{equation}
\mathcal{H}(t) = \bigl\{\,\mathrm{pk}_i \in \mathcal{L}(t) \;\mid\;
\mathrm{pk}_i \text{ was generated under ECDSA P-256}\,\bigr\},
\label{eq:hndl_set}
\end{equation}

and because the ledger is append-only,
$|\mathcal{H}(t)|$ is non-decreasing: $|\mathcal{H}(t_1)| \leq |\mathcal{H}(t_2)|$
for all $t_1 \leq t_2 \leq t_m$. Migration at $t_m$ prevents new
additions to $\mathcal{H}$ but does not reduce $|\mathcal{H}(t_m)|$,
which remains permanently positive. The blockchain-corrected security
parameter that should be substituted into equation~(\ref{eq:mosca}) is

\begin{equation}
Z_{\mathrm{blockchain}} = Z_{\mathrm{future}} +
(t_{\mathrm{eval}} - t_0) \geq Z_{\mathrm{future}},
\label{eq:mosca_blockchain}
\end{equation}

where $t_{\mathrm{eval}}$ is the date of the risk evaluation. For a
Fabric network operational since 2018 and evaluated in 2025,
$(t_{\mathrm{eval}} - t_0) = 7$~years of historical records are
permanently in $\mathcal{H}$, extending the effective $Z$ by 7~years
regardless of when migration is completed.

For a Fabric network operational since 2018, all ECDSA-signed records
from that date forward are permanently on the ledger as harvest material.
Even a perfect post-quantum migration completed today cannot protect
those records. The implication is that migration should be completed as
early as possible to minimise the volume of additional records committed
under classically vulnerable signatures, and that affected organisations
must accept some permanent residual exposure for historical records while
ensuring that all future records are protected.

\section{Formal Integration Requirements}
\label{sec:tm_requirements}

\subsection{Requirements Derivation}

The threat analysis, risk prioritisation, and timeline assessment jointly
establish a complete set of integration requirements for a quantum-resistant
Hyperledger Fabric deployment. Eight requirements are derived, each
traceable to one or more identified threat scenarios. Together they
constitute a complete and non-redundant specification: every threat
scenario from the risk matrix is addressed by at least one requirement,
and no requirement addresses a concern not traceable to an identified
threat.

\textbf{R1: Replace ECDSA P-256 with a NIST-standardised post-quantum
signature scheme at all signing points in the transaction lifecycle.}
R1 addresses the transaction signing vulnerabilities of TS2 and TS3.
Scope encompasses client proposal signing, peer endorsement signing,
orderer block signing, signature re-verification at the committing
peer, MSP channel-configuration signing, and gossip message signing,
corresponding to operations Op-01, Op-03, Op-05, Op-07, Op-09, Op-10,
Op-12, Op-14, Op-15, Op-22, and Op-23 in the inventory. The algorithm must be drawn from FIPS~204 (ML-DSA, finalised by NIST in
August 2024) or FIPS~206 (FN-DSA, issued as an NIST Initial Public Draft
in August 2025)~\cite{nist_fips204, nist_fips206}.
SLH-DSA from FIPS~205 is noted as a candidate for offline CA certificate
signing contexts where signing latency is not a concern, but is excluded
from the primary per-transaction requirement due to signature sizes of
7,856 to 49,856 bytes that would render Fabric transaction records
impractically large. The specific algorithm selection from among ML-DSA-44,
ML-DSA-65, and FN-DSA-512 is a deployment decision governed by the
algorithm selection framework of Chapter~\ref{ch:empirical}.

\textbf{R2: Replace ECDH P-256 with a NIST-standardised post-quantum key
encapsulation mechanism at all TLS key exchange points.}
R2 addresses the TLS confidentiality vulnerability of TS4 and the gossip
TLS component of TS5, corresponding to operations Op-02, Op-06, Op-08,
Op-11, Op-13, Op-20, and Op-24 in the inventory. The algorithm must be
ML-KEM-768 from FIPS~203, standardised by NIST in August
2024~\cite{nist_fips203}. During the transitional migration period, a
hybrid scheme combining classical ECDH with ML-KEM is required to maintain
interoperability with classical nodes. The hybrid approach ensures that
the TLS session key is derived from both the classical and post-quantum
key exchanges, meaning that an adversary must break both to recover the
session key. This provides protection against both classical and quantum
adversaries during the migration period. The \texttt{X25519MLKEM768}
hybrid key agreement scheme, standardised in TLS 1.3 by the IETF, is the
specified implementation of this hybrid approach~\cite{ietf_hybrid_tls}.

\textbf{R3: Issue X.509 certificates carrying post-quantum public keys
from the live Fabric CA.}
R3 specifically and exclusively addresses TS1. As established in the
partial migration vulnerability analysis, TS1 cannot be addressed by any
combination of BCCSP signing requirements; it requires that the CA itself
issue certificates under post-quantum keys at runtime. R3 must be
satisfied by a live CA capable of runtime certificate issuance, not by
an offline generation tool, because a migration that leaves the live CA
in ECDSA provides no protection for identities enrolled after the initial
network bootstrap.

\textbf{R4: Preserve the BCCSP plugin interface contract.}
R4 requires that all post-quantum algorithm integration be encapsulated
within a new BCCSP plugin that satisfies the existing BCCSP interface,
with no modification to the interface itself and no changes required to
components above the BCCSP layer. This architectural constraint reflects
the design intent of Fabric's BCCSP abstraction: the interface exists
precisely to decouple the cryptographic implementation from the components
that use it. Violating R4 by modifying the interface or introducing
post-quantum awareness into components above the BCCSP layer would create
a coupling between the cryptographic migration and the broader Fabric
upgrade process, increasing integration complexity and the risk of
introducing new vulnerabilities or functional regressions. R4 ensures the
integration is minimal, non-invasive, and compatible with the Fabric
component ecosystem without requiring changes to the Fabric SDK, the
chaincode framework, or the gossip protocol implementation.

\textbf{R5: Provide crypto-agility via configuration.}
R5 requires that no specific post-quantum algorithm be hardcoded. The
active algorithm must be selectable via configuration change without code
modification, enabling response to future algorithm deprecations, security
discoveries, or regulatory changes without requiring a software release.

\textbf{R6: Maintain functional correctness.}
R6 requires that post-quantum signed transactions follow the same logical
flow through the Fabric transaction lifecycle as classical transactions.
Endorsement policies, chaincode execution, ledger commitment, and MSP
identity validation must all function correctly with post-quantum
algorithm material.

\textbf{R7: Implement constant-time post-quantum cryptographic operations.}
R7 requires that all post-quantum signing and key generation operations
use constant-time algorithm implementations. This requirement is
particularly critical for FN-DSA, whose signing procedure involves
discrete Gaussian sampling over NTRU lattices, a procedure historically
associated with timing variability that can expose private key
material~\cite{pessl2019}. Compliance is achieved through careful library
selection and verification of constant-time properties.

\textbf{R8: Provide a documented phased migration pathway.}
R8 requires that the architecture define an incremental migration strategy
that does not require simultaneous upgrade of all Fabric components.
Enterprise Fabric deployments span multiple organisations whose upgrades
cannot be fully synchronised, and production networks cannot typically
sustain extended downtime for complete replacement. The three-phase design
that satisfies this requirement is specified in Chapter~\ref{ch:architecture}.

\subsection{Requirements Traceability}

Table~\ref{tab:req_traceability} presents the requirements traceability
matrix, mapping each requirement to the threat scenarios it addresses.
Following the table is a discussion of the key traceability insights
it reveals.

\begin{table}[ht]
\centering
\caption{Requirements traceability matrix for the eight formal integration
requirements derived from the quantum threat analysis.}
\label{tab:req_traceability}
\setlength{\tabcolsep}{5pt}
\footnotesize
\begin{tabular}{@{}cl>{\raggedright\arraybackslash}p{2.4cm}l@{}}
\toprule
\textbf{Req.} & \textbf{Description} & \textbf{Addresses} &
\textbf{Type} \\
\midrule
R1 & PQC signing at all signing points   & TS2, TS3          & Functional \\
R2 & PQC key encapsulation in TLS        & TS4, TS5          & Functional \\
R3 & PQC certificates from live CA       & TS1               & Functional \\
R4 & Preserve BCCSP interface            & Quality constraint & Architectural \\
R5 & Crypto-agility via configuration    & Quality constraint & Architectural \\
R6 & Functional correctness              & Quality constraint & Quality \\
R7 & Constant-time operations            & Side-channel       & Security \\
R8 & Phased migration pathway            & Operational        & Operational \\
\bottomrule
\end{tabular}
\end{table}

Two properties of the requirement set are confirmed by
Table~\ref{tab:req_traceability}. First, every threat scenario in the
risk matrix is addressed by at least one functional requirement: TS1 by
R3, TS2 and TS3 by R1, and TS4 and TS5 by R2. The requirement set is
therefore complete with respect to the identified threats. Second, the
traceability of R3 to TS1 alone confirms the partial migration
vulnerability: no subset of the requirements that omits R3 can claim
to address the full threat model. This is the formal basis for the
complete four-layer integration design of Chapter~\ref{ch:architecture},
in which R3 is satisfied by extending the live Fabric CA to issue
post-quantum certificates.

The five non-functional requirements, R4 through R8, are necessary for
a deployable, maintainable, and secure implementation. Without R4, the
integration risks breaking compatibility with components above the BCCSP
layer. Without R5, the implementation cannot respond to algorithmic
evolution. Without R7, the implementation may introduce new side-channel
vulnerabilities while closing the quantum vulnerability.

\section{Discussion}
\label{sec:tm_discussion}

\subsection{Completeness of the Threat Model}

The twenty-five-operation inventory produced in this chapter captures an
attack surface more comprehensive than any prior description in the
literature. Prior discussions of Fabric's quantum vulnerability have
focused on the transaction signing layer, the CA certificate layer, or
both, without systematic enumeration. This audit additionally identifies
the Raft inter-orderer consensus communications, the gossip protocol
signing and encryption, and the MSP channel configuration signing as
significant quantum attack surfaces, and provides component-level severity
and exploitability assessments for each.

The completeness of the inventory is verifiable by construction: the audit
enumerates every component in the standard Fabric transaction lifecycle,
maps each component's cryptographic operations, and classifies each
operation. No prior work has applied this systematic approach to Fabric
or any other enterprise blockchain platform. Campbell~\cite{campbell2019}
identified the general quantum threat to Fabric PKI but provided no
operational enumeration. Holcomb et al.~\cite{holcomb2021}, Hu~\cite{hu2025prototype},
and Abbas et al.~\cite{abbas2025} proceeded to implementation without
any threat model.

The significance of the completeness finding extends beyond academic
novelty. It has a direct practical implication: any integration proposal
that does not address all five threat scenario categories can be
evaluated against the inventory to identify which operations remain
vulnerable. For example, the proposal of Holcomb et al.\ addresses Op-05
at the BCCSP signing layer and partially addresses Op-14 and Op-15 at
the committing peer validation layer, but does not address Op-17 through
Op-22 at the CA and MSP layers, leaving TS1 entirely open. The proposal
of Abbas et al.\ addresses certificate issuance at Op-19 through offline
certificate generation but does not address the live CA signing at Op-17
and Op-18 for runtime enrolments, leaving TS1 partially open. The
inventory and threat scenario framework therefore serve as an evaluation
tool for existing and future proposals, enabling objective comparison
of their scope and completeness.

The uniform ECDLP dependence identified by the audit is a finding with
its own architectural significance. The fact that all twenty-five
operations use the same cryptographic primitive family means that Shor's
algorithm represents a single key that unlocks the entire security
architecture. By contrast, a cryptographically diverse architecture
where different layers use different algorithm families would require a
quantum adversary to develop separate attacks on each family. The
uniformity of Fabric's cryptographic architecture simplifies the
migration task but also means that no algorithm diversity provides any
partial protection: a network that has not migrated any component is
as vulnerable as a network that has migrated some components but not
all, because the unmigrated components share the same mathematical
vulnerability.

\subsection{The Ledger Immutability Amplification}

The most distinctive finding of this threat model relative to threat
models for conventional cryptographic systems is the permanent
amplification of the harvest-now-decrypt-later threat by the blockchain's
immutable ledger structure. In conventional systems, the HNDL threat is
mitigated through a combination of timely migration and key rotation. In
Fabric, no such mitigation is available for historical records. Every
ECDSA public key that has appeared on the ledger, including those of
participants who have been removed from the network, remains permanently
accessible as a harvest target.

This finding has a direct operational implication. The appropriate
response to the HNDL threat in a Fabric deployment is not to treat it
as a future problem that migration will eventually solve but to treat it
as a present problem that grows in severity with every additional block
committed under classical signatures. Migration protects future records;
it cannot retroactively protect historical ones.

This analysis leads to a counterintuitive governance implication. In
conventional PKI deployments, regulators and risk managers typically
accept some tolerance window between the completion of a migration and
the date on which the CRQC is projected to arrive, because classical
keys that have been retired cannot be targeted retrospectively. In Fabric
deployments, this tolerance window is zero for the historical portion of
the ledger, because the public keys are permanent. Risk managers for
Fabric deployments should therefore treat the effective HNDL harvest
window as open from the genesis block date, not from the current date,
when evaluating the regulatory and liability implications of delayed
migration.

The practical distinction between Fabric and conventional deployments
can be stated precisely. In a conventional TLS deployment, a key that
was used in 2020 and rotated in 2022 is no longer accessible as a
harvest target after 2022, because the public key was not published
to a permanent record. In a Fabric deployment, a signing key used in
2020 has its corresponding public key permanently embedded in every
transaction it signed, all of which remain on the ledger indefinitely.
The harvest window for that key closes only when the key itself is
replaced in future transactions, but the historical transactions signed
with the old key remain permanently vulnerable. For any Fabric deployment
carrying sensitive data, this distinction has significant implications
for the calculation of regulatory exposure and litigation risk.

\subsection{Implications for Integration Requirements}

The requirements derived from the threat analysis reflect a tighter
integration scope than any prior proposal. Specifically, R3 requires a
live CA capable of runtime post-quantum certificate issuance, not merely
the bootstrap-time certificate generation tools used by Holcomb et al.\
and Abbas et al. This requirement arises directly from the partial
migration vulnerability analysis: a deployment that relies on offline
certificate generation cannot issue post-quantum certificates to newly
enrolled identities and therefore leaves TS1 open for the certificates
of those new identities regardless of the state of the BCCSP signing
layer.

\subsection{Generalisability to Other Permissioned Blockchain Platforms}

While the specific operations enumerated in this audit are specific to
Hyperledger Fabric v2.5.12, the methodology, threat scenarios, and
requirements framework are designed to generalise across the family of
permissioned enterprise blockchain platforms. The core vulnerability
identified, the uniform dependence of certificate-based identity
management and transaction signing on the elliptic curve discrete
logarithm problem, is not specific to Fabric. Hyperledger Besu, the
Ethereum-compatible permissioned platform, uses ECDSA for transaction
signing and TLS for inter-node communications in the same way. Quorum
and Corda similarly rely on ECDSA-based certificate infrastructure.

The HNDL amplification analysis applies with equal force to all
blockchain platforms with append-only ledger structures, which is
the defining characteristic of all distributed ledger technologies.
Any platform that permanently preserves public keys and signatures
on an immutable ledger is subject to the same amplified harvest-now-
decrypt-later risk identified for Fabric.

The five threat scenario categories, identity impersonation via certificate
forgery, transaction signature forgery, block or transaction integrity
compromise via node key recovery, TLS session decryption, and peer
communication protocol manipulation, will recur in any permissioned
blockchain threat model, with platform-specific variations in the
specific operations affected and the severity of each scenario depending
on architectural details. The eight integration requirements derived for
Fabric, particularly the separation between certificate infrastructure
migration and signing layer migration, and the requirement for a live CA
capable of runtime post-quantum certificate issuance, are likely to be
relevant to migration planning for any certificate-based permissioned
blockchain platform.

\subsection{Relationship Between the Threat Model and Algorithm Selection}

The threat model developed in this chapter informs not only the
requirements for integration but also the algorithm selection decisions
described in Chapter~\ref{ch:architecture}. The severity assessment of
Op-14 and the analysis of signature verification at scale directly
motivate the preference for algorithms with smaller output sizes over
algorithms with faster signing operations. The critical requirement for
constant-time operations in R7 directly motivates the library selection
decisions for FN-DSA implementation. The requirement for crypto-agility
in R5 is driven by the recognition that the post-quantum standardisation
process may produce updates or revisions to algorithm specifications, as
occurred with earlier candidates during the NIST evaluation rounds.

The risk prioritisation matrix, which places TS1 and TS3 at joint P1
priority, explains why the three-phase migration design in
Chapter~\ref{ch:architecture} addresses the certificate infrastructure
in Phase 2 rather than deferring it to a later phase. A migration design
that addressed BCCSP signing first and certificate infrastructure only
later would leave the P1 threats unaddressed during the Phase 1 period.
The phased design ensures that the highest-priority threats are fully
addressed as early as possible in the migration sequence, with Phase 1
providing partial protection and Phase 2 completing the addressing of
the P1 threats.

\section{Conclusion}
\label{sec:tm_conclusion}

This chapter has presented the first component-level quantum vulnerability
audit and formal threat model for Hyperledger Fabric and for enterprise
blockchain platforms in general. The audit of Fabric v2.5.12 identifies
twenty-five quantum-vulnerable cryptographic operations distributed across
five transaction lifecycle phases, two identity management subsystems,
the gossip protocol, and the key management layer. The audit confirms
that Fabric's quantum vulnerability is complete and uniform: every
public-key cryptographic operation relies on the elliptic curve discrete
logarithm problem and is simultaneously threatened by Shor's algorithm.

The analysis identifies Op-10, the orderer block header signing operation,
and Op-17, the root CA certificate self-signing operation, as the
operations of highest strategic importance. Both are accessible to any
passive ledger observer, both have network-wide impact when compromised,
and both are subject to an unbounded harvest-now-decrypt-later window
that begins from network genesis and cannot be closed by post-migration
key rotation. The signature verification at scale analysis additionally
establishes that Op-14, the endorsement signature re-verification at
committing peers, is the operation of highest performance significance:
its throughput sensitivity to signature output size motivates the
performance analysis design of Chapter~\ref{ch:empirical}.

Five formal threat scenarios are constructed and assessed in a risk
prioritisation matrix. TS1 and TS3 share the joint P1 priority, both
rated Critical severity and High exploitability. The analysis of Mosca's
inequality under two CRQC timeline estimates confirms that healthcare,
financial settlement, and government procurement deployments require
immediate migration action regardless of which timeline materialises.
The harvest-now-decrypt-later threat is uniquely amplified by the
immutable ledger structure, and the correct formulation of Mosca's
inequality for blockchain deployments substitutes the maximum historical
sensitivity lifetime for the desired future sensitivity lifetime.

Eight formal integration requirements are derived and validated through
a requirements traceability matrix. The critical finding is that TS1
is addressed exclusively by R3, the requirement for a live Fabric CA
capable of runtime post-quantum certificate issuance. No subset of the
requirements that omits R3 can claim to address the full quantum attack
surface. This result formally justifies the complete four-layer
integration scope of Chapter~\ref{ch:architecture} and the decision
to extend both Hyperledger Fabric and Fabric CA source code to support
post-quantum algorithms at runtime. The threat taxonomy, risk
prioritisation, timeline analysis, and eight integration requirements
collectively constitute Contribution~C1 of this thesis.



\bibliographystyle{IEEEtran}
\bibliography{Ref}



\end{document}
